\clearpage
\section{Introduction}
\label{iii:sec:introduccion}

Une relation constitutive organise la réponse d’un système, mais sa
forme fonctionnelle et ses coefficients posent des questions différentes. Une
fois les coefficients connus, l’équation permet de calculer des réponses ;
déterminer pourquoi ces coefficients ont conjointement leurs valeurs
nécessite de décrire la structure qui les fixe. Ce travail aborde cette
seconde question au moyen d’une construction mathématique à deux canaux
orientés. Sa thèse est que la structure précédente détermine
l’équation et que l’évaluation conserve une partie récupérable de cette
structure. Le vide électromagnétique constitue ici une réalisation
précise de cette thèse : la permittivité, la perméabilité, l’impédance et la vitesse
s’obtiennent comme quatre lectures reliées d’une même réponse.

Le résultat principal comprend tant la construction directe que
l’inverse. Dans le domaine contractant, l’opérateur constitutif est positif,
ses valeurs propres appartiennent à un intervalle déterminé et ses quatre
coordonnées satisfont des identités exactes. La stricte monotonie de la
fonction de réponse fournit un inverse unique ; avec les feuillets
étiquetés, cet inverse retrouve les canaux et la paire angulaire. Cette
récupération permet de composer trois descriptions : la réponse du
vide, la forme de l’ellipse d’action et la fonctionnelle de
Barbero--Immirzi. La comparaison ne se limite donc pas à des valeurs
numériques : elle identifie les opérations qui transmettent l’orientation et
l’information spectrale entre ces descriptions.

\subsection{Noyau générateur et provenance de la réponse}

Le modèle se construit à partir d’un noyau formel mathématique appelé
holographie modulaire triadique, ci-après HMT. Son premier objet,
\emph{All Possible Paths} (APP), réunit des chemins et des évaluations
arithmétiques sur les $81$ positions du support
$X_9=(\mathbb Z/9\mathbb Z)^2$. Les translations de générateurs
$\{(1,0),(-1,0),(0,1),(0,-1)\}$ définissent son graphe de Cayley additif,
un réseau cyclique bidimensionnel de réalisation toroïdale. Les
évaluations somme et produit agissent sur ce support ; chacune conserve
son résidu et son quotient. La structure des chemins et les évaluations
ont ainsi des fonctions distinctes au sein d’APP : les déplacements
spécifient les chemins, tandis que les feuillets arithmétiques déterminent les
données qu’ils transportent.

Le deuxième objet, TRIT, est la signature ternaire orientée
$\{-1,0,+1\}$, obtenue au moyen de la représentation équilibrée des
classes modulo $3$. Sa réalisation quadratique distingue les types
hyperbolique, parabolique et elliptique, respectivement. Le troisième objet
est le \emph{Topological Phase Kernel}, ci-après TPK : il compose la
sélection de phase, le transport cardinal et la mise à jour de mémoire.
La dynamique agit sur les positions, les feuillets, les quotients, l’orientation,
les préfixes et les données de frontière. Le retour d’une phase conserve
l’histoire accumulée, de sorte qu’un tour peut fermer une coordonnée
sans restituer l’état complet.

Le prolongement nonadique transmet cette structure au moyen de préfixes et de
frontières compatibles. Dans son développement conjoint sont engendrées les
coordonnées $\pi_{\rm HMT}$, $\varphi_{\rm HMT}$,
$e_{\rm HMT}$ et $\alpha_{\rm HMT}$. Le registre dodécaphasique $K$
et son évaluation périodique interviennent dans la détermination de
$\alpha$ au moyen de la clôture dodécaphasique ; la représentation analytique complète exprime cette même
coordonnée dans la classe de queues construite. On obtient ensuite
les sections d’action et la paire angulaire. L’utilisation de ces résultats
comme antécédents conserve leurs démonstrations au sein de l’article.
Cette séquence fixe la généalogie des canaux avant d’évaluer la
réponse constitutive.

L’orientation des feuillets et leur incidence aréale--volumétrique remplissent
des rôles complémentaires. La première distingue le retour d’un chemin
de la restauration de son orientation ; la seconde provient du calendrier
des modes et de sa matrice d’incidence. La fréquence chronologique, la
mesure du système symbolique et la marque régionale de frontière maintiennent
leurs définitions propres. Ces distinctions, développées dans la
section~\ref{iii:sec:hojas}, précisent l’information qui parvient au
transport angulaire et évitent de la remplacer par un rapport scalaire isolé.

\clearpage
\subsection{Détermination constitutive et récupération de la forme}

Avec $T=q_+P_++q_-P_-$ et $0<q_\pm<1$, les incidences temporelles
$90$ et $120$ déterminent
\[
\mathcal V=(I-T^{90})(I-T^{120})^{-1}.
\]
En écrivant $S=T^{30}$, l’équation
$\mathcal V(I+S+S^2+S^3)=I+S+S^2$ montre la règle qui fixe la
réponse pour ce transport. La section~\ref{iii:sec:constitutiva}
démontre son existence et son unicité et établit la fonction scalaire
$f:(0,1)\to(3/4,1)$. Ses valeurs $r_\pm=f(q_\pm^{30})$
produisent
\[
\widehat\varepsilon=r_+^2,\qquad
\widehat\mu=r_-^2,\qquad
\widehat Z=r_-/r_+,\qquad
\widehat c=(r_+r_-)^{-1}.
\]
La positivité et les lectures de produit et de quotient déterminent les deux
carrés une fois les relations constitutives fixées. L’unicité
affirmée appartient à ce système explicite de règles et à son domaine.

L’inversion cubique retrouve $q_\pm$ puis les coordonnées
$x=-\frac12\log(q_+q_-)$,
$y=\frac12\log(q_-/q_+)$. Avec les étiquettes de feuillet conservées,
l’impédance et la vitesse suffisent à cette reconstruction.
La proposition~\ref{iii:prop:forma-LC} obtient ensuite
l’anisotropie de l’ellipse d’action et l’impédance inductive--capacitive
relative. Les références d’action, de fréquence et d’impédance complètent
sa réalisation dimensionnelle. La proposition~\ref{iii:prop:identidad-velocidad}
prouve, dans la même coordinatisation, $c_{\rm int}=c_{\rm vac}$ et relie la
longueur de trajectoire à la durée élémentaire. L’inversion
reconstruit ainsi des antécédents spectraux déterminés, tandis que le chemin
et la mémoire restent dans l’état enrichi dont ils proviennent.

Cet inverse induit également une loi de composition entre réponses :
la multiplication des transports s’exprime au moyen de l’opération
transportée et admet une coordonnée additive. La
section~\ref{amp-higgs-seccion} compose ensuite les mêmes canaux avec
les signatures angulaires de jauge et de Higgs. L’inverse conjoint, son
domaine exact et la reconstruction de l’impédance conservent les conditions
d’échelle, de raffinement et de normalisation de cette réalisation.

\subsection{Cycle, constantes et fonctionnelles orientées}

La fonction $f$ admet une réalisation sur le cycle à douze phases.
Ses sous-espaces de dimensions $3$ et $4$ produisent le quotient de
déterminants ; leur différence de trace normalisée donne $-1/12$.
Dans un plan orienté du même cycle, un quart de tour $J$ satisfait
$J^2=-I$ et sa résolvante détermine $s/(1+s^2)$. Deux intégrations
logarithmiques, avec leurs conditions à l’origine, engendrent le moment
$\mathcal C_2(s)$ et sa valeur à l’extrémité, la constante de Catalan.
Le théorème~\ref{iii:thm:catalan-pell} relie exactement
l’évaluation en $2-\sqrt3$ à cette valeur à l’extrémité et à
$\log(2+\sqrt3)$. La section algébrique de Pell et les arguments
angulaires $q_\pm^{30}$ sont des évaluations distinctes de la même famille.

Les traces des puissances du transport déterminent en outre un tableau
périodique qui reconstruit le logarithme du quotient constitutif.
Sa série de moments, la représentation de Mellin et les bornes du reste
sont démontrées dans la section~\ref{amp-afin-sec}. La composition diagonale
et la tensorielle maintiennent explicite la profondeur conservée par chaque
opération ; le retour spectral demeure distinct du retour de l’état.

Le théorème~\ref{iii:thm:restitucion-funcional-catalan} complète la
relation au moyen d’une restitution intégrale : la fonction constitutive
fixe le moment orienté entier et la condition de limite nulle à
l’origine détermine sa normalisation. La preuve distingue la récupération
de cette famille de l’inversion de ses deux évaluations de canal.
Les deux opérations conservent la lecture de vitesse
$\widehat c=(\det\mathcal V)^{-1}$ et sa trace normalisée, qui relient
le transport constitutif à la durée et à la longueur du même état.

La composition avec Barbero--Immirzi conserve une information supplémentaire :
le signe relatif des canaux. L’incidence hexade--octade et le
calendrier fixent les termes de la fonctionnelle angulaire ; l’inverse de
$f$ permet de l’exprimer sur $\mathcal V$. Le
théorème~\ref{iii:thm:barbero-factorization} démontre cette
factorisation. Un échange de canaux par rapport à une marque
d’orientation fixe inverse le signe de la fonctionnelle ; une conjugaison
simultanée de l’opérateur et de la marque le conserve. De cette manière, la réponse
constitutive transporte l’information nécessaire à sa lecture angulaire.

% BEGIN SINTESIS_ADITIVA_20260922
\par
La section~\ref{delta22:iii:restitucion} réunit le potentiel spectral de la réponse constitutive, la restitution globale des canaux au moyen de la vitesse normalisée et de Barbero--Immirzi, ainsi que ses bornes de stabilité. Sur l’image engendrée, avec les étiquettes et les bases conservées, la composition retrouve la paire angulaire, la section d’action et la lecture périodique du registre dodécaphasique.
% END SINTESIS_ADITIVA_20260922

La section~\ref{iii:sec:incidencias-espectrales} rend explicites les
normalisations de cette composition. La factorisation positive sépare
le facteur commun de la réponse et sa composante unimodulaire orientée ;
les deux données permettent de restituer les quatre coordonnées constitutives.
La représentation de Fourier distingue ensuite les modes des
incidences $90/120$, avec les poids signés de la chaîne et une
normalisation de transformée déclarée. Cette organisation permet de
comparer les réalisations antécédentes à l’opérateur utilisé
dans l’article, en conservant pour chacune sa fonction et son domaine.

L’appendice~\ref{iii:hist:antecedentes} développe en outre la réalisation
finie sur les positions, les modes et la phase : il prouve son équivariance et la
décomposition en vingt-sept composantes, conserve les phases de ses
valeurs propres et évalue sa sensibilité paramétrique. Les lecteurs angulaires
antécédents sont reproduits avec leurs opérations et données déclarées.
Leur comparaison aux lecteurs constitutifs du corps du texte distingue les
poids, les normalisations et les domaines de chaque réalisation.

\subsection{Polarisation, charge et réalisation dimensionnelle}

Le développement électrique s’appuie également sur deux résultats des
feuillets arithmétiques. Le décalage entre capacités nonadique et décimale
produit une discordance cumulée uniformément bornée ; sa différence
temporelle définit l’accroissement de polarisation. D’autre part,
l’appariement des fibres additive et multiplicative dans $D_9^3$
détermine le mode tritique et la norme $\sqrt3/4$ du commutateur de leurs
projections. La section~\ref{iii:sec:vacancias-prefactor} conserve
les preuves et les domaines des deux constructions. Leur relation avec
la limite $3/4$ de la réponse temporelle précise la donnée qu’elles partagent
sans identifier leurs espaces opératoriels.

La section~\ref{iii:sec:carga} compose l’impédance, $\alpha$ et la
constante de Planck au moyen de l’unique section positive satisfaisant
$Z_{\rm vac}e_a^2=2\alpha h_a$. On en obtient la résistance de
von Klitzing, le quantum de conductance, le flux magnétique et la constante de
Josephson, avec des identités croisées exactes. Le retour de l’action
laisse invariants les deux premiers et transforme les deux derniers par
des puissances opposées du quotient d’action. Ce contenu électrique
est une conséquence des sections précédentes, non une liste de
constantes indépendantes.

L’évaluation finale maintient les droites d’action, de charge et de vitesse,
ainsi que les transformations de leurs coordonnées lors du changement d’unités.
Le tableau de la section~\ref{iii:sec:evaluacion} présente les coordonnées
internes obtenues par évaluation. Une comparaison physique se formule ensuite, avec la coordinatisation
dimensionnelle et le régime de couplage déclarés. L’argument
s’organise autour d’une question commune : quelle structure détermine la
réponse, quelles équations elle impose de satisfaire et quelle information peut
être retrouvée à partir des valeurs ainsi obtenues.
